\begin{theorem}[Two-label structure coefficients whose length dependence survives positive rescaling]%
    \label{thm:mpdo_twisted_dimer_two_label_rescaling_stable}
    \lean{MPOTensor.TwistedDimer.IsSameChannel}
    \lean{MPOTensor.TwistedDimer.alpha}
    \lean{MPOTensor.TwistedDimer.beta}
    \lean{MPOTensor.TwistedDimer.twoLabelChi}
    \lean{MPOTensor.TwistedDimer.twoLabelChi_posEntries}
    \lean{MPOTensor.TwistedDimer.twoLabelCoeffs}
    \lean{MPOTensor.TwistedDimer.twoLabelCoeffs_coeff}
    \lean{MPOTensor.TwistedDimer.twoLabelCoeffs_not_lengthIndependent}
    \lean{MPOTensor.TwistedDimer.twoLabelChiScaled}
    \lean{MPOTensor.TwistedDimer.rescaledCoeffs}
    \lean{MPOTensor.TwistedDimer.rescaledCoeffs_coeff}
    \lean{MPOTensor.TwistedDimer.%
        twoLabelCoeffs_rescaling_stable_not_lengthIndependent}
    \leanok
    \uses{def:mpdo_bnt_label_coefficient_family, def:mpdo_bnt_label_length_independent}
    Let the labels be $\{0,1\}$, with addition modulo two, and define
    positive diagonal matrices of size one, identified with their entries, by
    \begin{align}
        \chi_{ff',\,f+f'}   & =\alpha=\frac{7}{10},
        \notag                                      \\
        \chi_{ff',\,f+f'+1} & =\beta=\frac1{10}.
        \notag
    \end{align}
    Then $c^{(L)}_{ff'g}=\tr(\chi_{ff'g}^L)$ equals $\alpha^L$ when
    $g=f+f'$ and $\beta^L$ otherwise. These coefficients are not length
    independent. For every pair of positive rescalings $(s_0,s_1)$,
    replacing $\chi_{ff'g}$ by $(s_fs_{f'}/s_g)\chi_{ff'g}$ still gives
    coefficients that are not length independent. Unlike the
    one-label families of
    Theorems~\ref{thm:mpdo_length_dependent_rfp_example} and
    \ref{thm:mpdo_bnt_label_rescaling_stable_length_dependence}, the
    obstruction here is not a diagonal matrix with two distinct entries but
    the inconsistency of the equations
    $s_fs_{f'}\chi_{ff'g}=s_g$ over the eight triples of labels.
    This statement concerns the coefficients alone.
    Theorem~\ref{thm:mpdo_twisted_dimer_flag_sector_product_law} identifies
    them as the structure coefficients for the closed operators of the flag
    sectors of the twisted quantum dimer.
\end{theorem}
\begin{proof}\leanok
    \uses{def:mpdo_bnt_label_coefficient_family, def:mpdo_bnt_label_length_independent}
    If the powers of a positive real number $t$ agree at lengths one and
    two, then $t=t^2$, hence $t=1$. The triples $(0,0,0)$ and
    $(0,1,0)$ give $s_0\alpha=1$ and $s_1\beta=1$; the triple $(0,0,1)$ then
    gives $s_0^2\beta/s_1=\beta^2/\alpha^2=1$, contradicting
    $\alpha\ne\beta$.
\end{proof}

\begin{definition}[Flag sectors of the twisted quantum dimer]%
    \label{def:mpdo_twisted_dimer_flag_sectors}
    \lean{MPOTensor.TwistedDimer.mu}
    \lean{MPOTensor.TwistedDimer.flagWeight}
    \lean{MPOTensor.TwistedDimer.flagWeight_physIdx}
    \lean{MPOTensor.TwistedDimer.flagCoef}
    \lean{MPOTensor.TwistedDimer.unitTensor}
    \lean{MPOTensor.TwistedDimer.block_eq_unitTensor}
    \lean{MPOTensor.TwistedDimer.flagMPO}
    \lean{MPOTensor.TwistedDimer.flagMPO_apply_eq_T}
    \lean{MPOTensor.TwistedDimer.flagFamily}
    \lean{MPOTensor.TwistedDimer.verticalBNTMPO_flagFamily}
    \lean{MPOTensor.TwistedDimer.flagOperatorFamily}
    \lean{MPOTensor.TwistedDimer.flagOperatorFamily_operator}
    \leanok
    \uses{def:mpdo_twisted_dimer_tensor, def:mpdo_bnt_label_operator_family}
    The matrices of the twisted quantum dimer of
    Definition~\ref{def:mpdo_twisted_dimer_tensor} are diagonal in the flag
    qubit. Viewed vertically, the tensor therefore splits along the flag
    value $f\in\{0,1\}$ into two tensors $M_f$ acting on the two-qubit
    space of $L$ and $R$, with matrices indexed by pairs of horizontal bond
    indices. Write $\mu=\frac58$, so that $\mu^2=(x^2+y^2)/2$ at
    $x=\frac78$, $y=\frac18$. The rescaled sector $\widehat M_f=M_f/\mu$
    has, at the bond pair $((p,p',k),(q,q',k'))$, the matrix
    \begin{align}
        \widehat M_f^{(p,p',k),(q,q',k')}
         & =[k=k']\,\frac1{2\mu}\,C_k[p,p']\,(\tau_k)_{ff}
        E_{pp'}\otimes E_{qq'},
        \notag
    \end{align}
    where $[k=k']$ is one when $k=k'$ and zero otherwise.
    The matrix unit has row $(p,q)$ and column $(p',q')$.
    The entry of this matrix between the qubit configurations $(l,r)$ and
    $(l',r')$ is $\mu^{-1}$ times the bond-matrix entry of the twisted
    dimer at the physical indices $((l,r,f),(l',r',f))$ carrying the common
    flag value $f$. The horizontal blocks of
    Definition~\ref{def:mpdo_twisted_dimer_tensor} and the two vertical flag
    sectors are all matrix-unit tensors of this form, differing only in their
    one-site coefficients. The closed operators $O_L(\widehat M_f)$ on a ring
    of $L$ sites give the operators indexed by these two labels in
    Definition~\ref{def:mpdo_bnt_label_operator_family}. Their
    basis-of-normal-tensors properties are proved below, rather than assumed in
    this definition. This tensor is motivated by the question of
    \cite[lines~995--1010]{Cirac2016MPDO_arXiv}.
\end{definition}

\begin{lemma}[Cyclic path expansion of a trace]%
    \label{lem:matrix_cyclic_path_sum}
    \lean{Matrix.sum_fin_cons}
    \lean{Matrix.ofFn_prod_mulVec_apply}
    \lean{Matrix.trace_ofFn_prod_eq_sum_cyclic}
    \leanok
    Let $M_0,\ldots,M_{L-1}$ be square matrices of the same finite size over a
    commutative semiring, with $L\ge1$. Applying the ordered product
    $M_0\cdots M_{L-1}$ to a vector $v$ and reading the entry at $a$ sums,
    over all paths $(t_0,\ldots,t_L)$ with $t_0=a$, the product of the
    entries $(M_n)_{t_n t_{n+1}}$ along the path times $v_{t_L}$.
    Consequently,
    \begin{align}
        \tr(M_0\cdots M_{L-1})
         & =\sum_{t_0,\ldots,t_{L-1}}\prod_{n=0}^{L-1}(M_n)_{t_n t_{n+1}},
        \notag
    \end{align}
    with the cyclic convention $t_L=t_0$.
\end{lemma}
\begin{proof}\leanok
    The open expansion is an induction on $L$: splitting a path into its first
    step and the remaining path gives
    $(M_0(M_1\cdots M_{L-1})v)_a
        =\sum_c(M_0)_{ac}((M_1\cdots M_{L-1})v)_c$.
    For the trace, separate the last factor $M_{L-1}$ and apply the open
    expansion to the vector $c\mapsto(M_{L-1})_{ca}$, using the empty
    product when $L=1$. Splitting a cyclic assignment into its first entry
    $a$ and the remaining entries identifies the two sides term by term,
    with the last factor closing the cycle back to $a$.
\end{proof}

\begin{lemma}[Closed operators of the flag sectors]%
    \label{lem:mpdo_twisted_dimer_flag_sector_operators}
    \lean{MPOTensor.TwistedDimer.unitTensor_mul_single}
    \lean{MPOTensor.TwistedDimer.evalWord_unitTensor_ofFn}
    \lean{MPOTensor.TwistedDimer.mpo_unitTensor_apply}
    \lean{MPOTensor.TwistedDimer.prod_flagCoef}
    \lean{MPOTensor.TwistedDimer.mpo_flagMPO_apply}
    \leanok
    \uses{def:mpdo_twisted_dimer_flag_sectors, def:mpdo}
    For a matrix-unit tensor with one-site coefficients $c$, the ordered
    product along a nonempty pair of words vanishes unless both words
    satisfy the open bond-matching condition. When they do, the product is
    the matrix unit whose row is given by the left bits of the first indices
    and whose column is given by the right bits of the last indices,
    multiplied by the product of the coefficients. On a ring of positive
    length the closed operator therefore has entries
    $\Id[\text{cyclic bond match}](\sigma,\tau)
        \prod_n c(\sigma_n,\tau_n)$, where the first factor is the indicator
    that both words satisfy the cyclic bond-matching condition.
    In particular, for a positive length $L$ and words
    $\sigma_n=(p_n,p'_n,k_n)$, $\tau_n=(q_n,q'_n,k'_n)$ of
    horizontal bond indices,
    \begin{align}
        O_L(\widehat M_f)_{\sigma\tau}
         & =\Id[\text{cyclic bond match}](\sigma,\tau)
        \prod_n[k_n=k'_n]
        \prod_n\frac{C_{k_n}[p_n,p'_n]\,(\tau_{k_n})_{ff}}{2\mu}.
        \notag
    \end{align}
\end{lemma}
\begin{proof}\leanok
    \uses{def:mpdo_twisted_dimer_flag_sectors}
    A matrix unit multiplied on the left by a tensor matrix is a matrix unit
    or zero according to whether the right bits of that matrix match the row
    of the unit. The open-chain formula follows by induction on the length
    of the words, as for the twisted dimer itself in
    Theorem~\ref{thm:mpdo_twisted_dimer_ismpdo}. Taking the trace imposes the
    cyclic bond match. For a flag sector, the coefficient at a bond pair is
    the one-site weight of the first index when the block labels agree and
    zero otherwise. Thus the product of coefficients along the ring is the
    product of the weights of the first word if the two label strings agree,
    and zero otherwise.
\end{proof}

\begin{theorem}[Product law for the closed operators of the flag sectors]%
    \label{thm:mpdo_twisted_dimer_flag_sector_product_law}
    \lean{MPOTensor.TwistedDimer.tau_one_mul_self}
    \lean{MPOTensor.TwistedDimer.tau_add}
    \lean{MPOTensor.TwistedDimer.Cmat_eq_signs}
    \lean{MPOTensor.TwistedDimer.Cc}
    \lean{MPOTensor.TwistedDimer.ofFn_Cc_prod_apply}
    \lean{MPOTensor.TwistedDimer.trace_ofFn_Cc_prod}
    \lean{MPOTensor.TwistedDimer.sum_cyclic_Cmat}
    \lean{MPOTensor.TwistedDimer.prod_flagWeight_physIdx}
    \lean{MPOTensor.TwistedDimer.sum_sector}
    \lean{MPOTensor.TwistedDimer.flagWeight_add}
    \lean{MPOTensor.TwistedDimer.alpha_eq_x_div}
    \lean{MPOTensor.TwistedDimer.beta_eq_y_div}
    \lean{MPOTensor.TwistedDimer.mpo_flagMPO_mul}
    \lean{MPOTensor.TwistedDimer.flagOperatorFamily_hasSameLengthProductForm}
    \leanok
    \uses{def:mpdo_twisted_dimer_flag_sectors, def:mpdo_bnt_label_same_length_product_form, thm:mpdo_twisted_dimer_two_label_rescaling_stable}
    For every positive length $L$ and all flag values $f,f'\in\{0,1\}$, the
    closed operators of the two rescaled flag sectors satisfy
    \begin{align}
        O_L(\widehat M_f)\,O_L(\widehat M_{f'})
         & =\alpha^L\,O_L(\widehat M_{f+f'})
        +\beta^L\,O_L(\widehat M_{f+f'+1}),
        \notag
    \end{align}
    where $\alpha=\frac{7}{10}$, $\beta=\frac1{10}$, and the flag values are
    added modulo two. Thus the operators of
    Definition~\ref{def:mpdo_twisted_dimer_flag_sectors} satisfy the
    product law of
    Definition~\ref{def:mpdo_bnt_label_same_length_product_form}, in the form
    of \cite[Theorem~4.14(ii), lines~972--985]{Cirac2016MPDO_arXiv}, with the
    structure coefficients of
    Theorem~\ref{thm:mpdo_twisted_dimer_two_label_rescaling_stable}.
    These coefficients depend on $L$, and no positive rescaling of the two
    sectors makes them length independent. This product law alone does not
    identify the sectors as a basis of normal tensors; that property is
    established below.
\end{theorem}
\begin{proof}\leanok
    \uses{lem:mpdo_twisted_dimer_flag_sector_operators, lem:matrix_cyclic_path_sum, thm:mpdo_twisted_dimer_two_label_rescaling_stable}
    Fix cyclically bond-matched words $\sigma$, $\tau$ with the common block
    labels $k_n$ and left bits $p_n$ of $\sigma$; all other entries of both
    sides vanish by Lemma~\ref{lem:mpdo_twisted_dimer_flag_sector_operators}.
    The entry of the product is a sum over a middle word $\rho$, which must be
    cyclically bond matched with the same labels $k_n$, hence of the form
    $\rho_n=(t_n,t_{n+1},k_n)$ for a bit string $t$. Its contribution factors
    as the weight of $\sigma$ in the sector $f$ times
    $(2\mu)^{-L}\prod_n(\tau_{k_n})_{f'f'}\prod_n C_{k_n}[t_n,t_{n+1}]$.
    By Lemma~\ref{lem:matrix_cyclic_path_sum}, the sum over $t$ of the last
    product is $\tr(C_{k_0}\cdots C_{k_{L-1}})$. Since
    $C_k=xP_++(-1)^kyP_-$ with orthogonal rank-one idempotents $P_\pm$, the
    ordered product is $x^LP_++(-1)^{\sum_n k_n}y^LP_-$, with trace
    $x^L+(-1)^{\sum_n k_n}y^L$; this is verified entrywise through the sign
    expansion $C_k[p,q]=\frac12(x+(-1)^{k+p+q}y)$. The flag signs multiply
    as $(\tau_k)_{ff}(\tau_k)_{f'f'}=(\tau_k)_{gg}$ with $g=f+f'$, and the
    sign $(-1)^{\sum_n k_n}=\prod_n(\tau_{k_n})_{11}$ of the second summand
    shifts the label to $g+1$. With $x/(2\mu)=\frac{7}{10}$ and
    $y/(2\mu)=\frac1{10}$, the two summands are $\alpha^L$ and $\beta^L$ times
    the entries of $O_L(\widehat M_{f+f'})$ and $O_L(\widehat M_{f+f'+1})$.
    These are precisely the coefficients for $g=f+f'$ and $g=f+f'+1$ in
    Theorem~\ref{thm:mpdo_twisted_dimer_two_label_rescaling_stable}.
\end{proof}

\begin{theorem}[Horizontal canonical form of the twisted quantum dimer]%
    \label{thm:mpdo_twisted_dimer_horizontal_cf}
    \lean{MPOTensor.TwistedDimer.normalizedBlock}
    \lean{MPOTensor.TwistedDimer.normalizedBlock_apply}
    \lean{MPOTensor.TwistedDimer.toMPSTensor_apply_finProdFinEquiv}
    \lean{MPOTensor.TwistedDimer.coef_flag_zero_ne_zero}
    \lean{MPOTensor.TwistedDimer.block_toMPSTensor_isInjective}
    \lean{MPOTensor.TwistedDimer.normalizedBlock_toMPSTensor_isInjective}
    \lean{MPOTensor.TwistedDimer.star_coef}
    \lean{MPOTensor.TwistedDimer.conjTranspose_mul_self_normalizedBlock}
    \lean{MPOTensor.TwistedDimer.rightBitEquiv}
    \lean{MPOTensor.TwistedDimer.sum_single_eq_single_sum}
    \lean{MPOTensor.TwistedDimer.sum_normalizedCoef_sq}
    \lean{MPOTensor.TwistedDimer.normalizedBlock_isLeftCanonical}
    \lean{MPOTensor.TwistedDimer.normalizedBlock_isNormalTensor}
    \lean{MPOTensor.TwistedDimer.physTraceTransfer_normalizedBlock_one}
    \lean{MPOTensor.TwistedDimer.exists_eq_physIdx}
    \lean{MPOTensor.TwistedDimer.blockBondEquiv}
    \lean{MPOTensor.TwistedDimer.retainedBondEquiv}
    \lean{MPOTensor.TwistedDimer.ambientCoisometry}
    \lean{MPOTensor.TwistedDimer.conjTranspose_toMatrix_toPEquiv}
    \lean{MPOTensor.TwistedDimer.ambientCoisometry_mul_conjTranspose}
    \lean{MPOTensor.TwistedDimer.conjTranspose_mul_mul_ambientCoisometry}
    \lean{MPOTensor.TwistedDimer.finSigmaFinEquiv_symm_retainedBondEquiv_symm}
    \lean{MPOTensor.TwistedDimer.toMPSTensor_T_reconstruct}
    \lean{MPOTensor.TwistedDimer.canonicalFormData}
    \lean{MPOTensor.TwistedDimer.T_toMPSTensor_isCPSVCanonicalForm}
    \leanok
    \uses{def:mpdo_twisted_dimer_tensor, def:cpsv_canonical_form, def:nt_cpsv, def:left_canonical_tensor}
    Let $B_0$, $B_1$ denote the two horizontal blocks of the twisted quantum
    dimer $M$ of Definition~\ref{def:mpdo_twisted_dimer_tensor}, read as
    tensors with the doubled physical index $((l,r,f),(l',r',f'))$ and bond
    dimension four, and let $\widehat B_k=\frac85\,B_k$.
    The matrices of each $\widehat B_k$ span the full algebra of
    $4\times4$ matrices, and
    $\sum_{i,j}(\widehat B_k^{ij})^\dagger\widehat B_k^{ij}=\Id$.
    Hence each $\widehat B_k$ is a normal tensor with spectral radius one.
    The doubled-index tensor is in canonical form in the sense of
    Definition~\ref{def:cpsv_canonical_form}: with the permutation $U$ of the
    bond index $(p,p',k)$ into the block label $k$ followed by the block bond
    index $(p,p')$,
    \begin{align}
        M^{ij} & =U^\dagger\Bigl(\tfrac58\,\widehat B_0^{ij}\oplus
        \tfrac58\,\widehat B_1^{ij}\Bigr)U.
        \notag
    \end{align}
    The physical-trace transfer of $\widehat B_1$ vanishes.
\end{theorem}
\begin{proof}\leanok
    \uses{thm:mpdo_twisted_dimer_closures, thm:is_normal_tensor_of_is_normal_left_canonical}
    The matrix of $B_k$ at the physical pair $((p,q,0),(p',q',0))$ is the
    matrix unit at the bond pair $((p,p'),(q,q'))$ with the nonzero
    coefficient $\frac12 C_k[p,p']$, so these matrices span the bond algebra.
    The product $(\widehat B_k^{ij})^\dagger\widehat B_k^{ij}$ is the
    diagonal matrix unit at the right bits $(r,r')$ of $i$ and $j$, with
    coefficient $\frac{64}{25}$ times the squared coefficient of the tensor
    matrix. Grouping by the right bits, the coefficients at fixed
    $(r,r')$ sum to
    $\frac{64}{25}\cdot2\cdot\frac14(2C_k[0,0]^2+2C_k[0,1]^2)=1$,
    so the sum of the products is the identity.
    Theorem~\ref{thm:is_normal_tensor_of_is_normal_left_canonical} gives
    normality with spectral radius one. The reconstruction is the block
    decomposition of Definition~\ref{def:mpdo_twisted_dimer_tensor} entry by
    entry, with $\frac58\cdot\frac85=1$ restoring the original blocks.
    The matrix $U$ is a permutation matrix, hence a coisometry.
    The physical-trace transfer of $\widehat B_1$ is $\frac85$ times that of
    $B_1$, which vanishes by
    Theorem~\ref{thm:mpdo_twisted_dimer_closures}.
\end{proof}

\begin{lemma}[Blocking raises the physical-trace transfer to a power]%
    \label{lem:mpdo_phys_trace_transfer_blocking}
    \lean{MPOTensor.physTraceTransfer_blockTensor}
    \leanok
    \uses{def:mpdo_physical_blocking, def:mpo_physical_trace_transfer}
    For every tensor $M$ and every blocking length $L$, the physical-trace
    transfer of the $L$-site blocking of $M$ is the $L$-th power of the
    physical-trace transfer of $M$.
\end{lemma}
\begin{proof}\leanok
    \uses{lem:mpdo_trace_loop_chain}
    The physical-trace transfer of the blocking is the sum over all
    length-$L$ configurations of the diagonal word evaluations, which is the
    $L$-th power of the one-site loop by
    Lemma~\ref{lem:mpdo_trace_loop_chain}.
\end{proof}

\begin{theorem}[The twisted quantum dimer is not simple]%
    \label{thm:mpdo_twisted_dimer_not_simple}
    \lean{MPOTensor.TwistedDimer.blockedCanonicalFormData}
    \lean{MPOTensor.TwistedDimer.blockedCanonicalFormData_blocks}
    \lean{MPOTensor.TwistedDimer.%
        doubledPhysTraceTransfer_blockedCanonicalFormData_blocks_one}
    \lean{MPOTensor.TwistedDimer.T_not_isSimple}
    \leanok
    \uses{def:mpdo_simple_tensor, def:mpdo_twisted_dimer_tensor}
    The twisted quantum dimer of Definition~\ref{def:mpdo_twisted_dimer_tensor}
    is not simple in the sense of Definition~\ref{def:mpdo_simple_tensor}: for
    every positive blocking length $L$ and every basis-of-normal-tensors sector presentation of
    the doubled-index tensor of the $L$-site blocking, some representative
    has nilpotent physical-trace transfer.
\end{theorem}
\begin{proof}\leanok
    \uses{thm:mpdo_twisted_dimer_horizontal_cf, thm:cpsv_literal_canonical_form_positive_blocking, thm:cpsv_physical_reindexing, thm:mpdo_blocking_doubled_index, thm:cpsv_canonical_form_bnt_refinement, thm:cpsv_representative_sector_decomposition, thm:mpdo_bnt_nonnilpotency_independent, thm:mpdo_doubled_phys_trace_nilpotent_gauge_phase, lem:mpdo_doubled_phys_trace_transfer_to_mps, lem:mpdo_phys_trace_transfer_blocking}
    Fix $L>0$ and a presentation $P$. Blocking the canonical form of
    Theorem~\ref{thm:mpdo_twisted_dimer_horizontal_cf} by $L$ sites
    (Theorem~\ref{thm:cpsv_literal_canonical_form_positive_blocking}) and
    relabeling the doubled alphabet along the pairing of the blocked ket and
    bra words (Theorems~\ref{thm:cpsv_physical_reindexing}
    and~\ref{thm:mpdo_blocking_doubled_index}) gives a canonical form of the
    blocked doubled-index tensor whose retained blocks are the $L$-site
    blockings of $\widehat B_0$ and $\widehat B_1$. Its
    basis-of-normal-tensors refinement
    (Theorem~\ref{thm:cpsv_canonical_form_bnt_refinement}) is a
    basis-of-normal-tensors sector
    presentation $Q$
    (Theorem~\ref{thm:cpsv_representative_sector_decomposition}) in which
    every retained block is related to the representative of its phase class
    by a similarity transformation and a phase.
    The blocked $\widehat B_1$ has vanishing physical-trace
    transfer by Lemma~\ref{lem:mpdo_phys_trace_transfer_blocking} together
    with Lemma~\ref{lem:mpdo_doubled_phys_trace_transfer_to_mps}, so by
    Theorem~\ref{thm:mpdo_doubled_phys_trace_nilpotent_gauge_phase} the
    representative of its phase class has nilpotent physical-trace transfer.
    Theorem~\ref{thm:mpdo_bnt_nonnilpotency_independent} transfers this to
    $P$.
\end{proof}

\begin{theorem}[Vertical canonical form of the twisted quantum dimer]%
    \label{thm:mpdo_twisted_dimer_vertical_cf}
    \lean{MPOTensor.TwistedDimer.sectorDim}
    \lean{MPOTensor.TwistedDimer.sectorMult}
    \lean{MPOTensor.TwistedDimer.sectorWeight}
    \lean{MPOTensor.TwistedDimer.RetainedCoordinate}
    \lean{MPOTensor.TwistedDimer.sectorCoord}
    \lean{MPOTensor.TwistedDimer.sectorCoord_apply}
    \lean{MPOTensor.TwistedDimer.physCoordEquiv}
    \lean{MPOTensor.TwistedDimer.physCoordEquiv_sectorCoord}
    \lean{MPOTensor.TwistedDimer.verticalCoisometry}
    \lean{MPOTensor.TwistedDimer.verticalCoisometry_apply}
    \lean{MPOTensor.TwistedDimer.verticalCoisometry_mul_conjTranspose}
    \lean{MPOTensor.TwistedDimer.conjTranspose_mul_verticalCoisometry}
    \lean{MPOTensor.TwistedDimer.verticalCoisometry_conj}
    \lean{MPOTensor.TwistedDimer.T_apply_eq_zero_of_bitF_ne}
    \lean{MPOTensor.TwistedDimer.flagFamily_finProdFinEquiv}
    \lean{MPOTensor.TwistedDimer.verticalCoisometry_conj_verticalTensor}
    \lean{MPOTensor.TwistedDimer.verticalTensor_T_eq_conjTranspose_mul}
    \lean{MPOTensor.TwistedDimer.mpv_verticalTensor_T}
    \lean{MPOTensor.TwistedDimer.Cmat_ne_zero}
    \lean{MPOTensor.TwistedDimer.flagWeight_physIdx_zero_ne_zero}
    \lean{MPOTensor.TwistedDimer.exists_flagFamily_eq_smul_single}
    \lean{MPOTensor.TwistedDimer.flagFamily_isInjective}
    \lean{MPOTensor.TwistedDimer.flagFamily_isNormal}
    \lean{MPOTensor.TwistedDimer.flagWeight_physIdx_zero_zero_zero}
    \lean{MPOTensor.TwistedDimer.flagWeight_physIdx_zero_zero_one}
    \lean{MPOTensor.TwistedDimer.linearIndependent_mpvState_flagFamily}
    \lean{MPOTensor.TwistedDimer.flagFamily_isBNT}
    \lean{MPOTensor.TwistedDimer.T_isVerticalCF}
    \leanok
    \uses{def:mpdo_twisted_dimer_flag_sectors, def:vertical_cf_mpdo, def:vertical_tensor_mpdo, def:bnt}
    The twisted quantum dimer $M$ of
    Definition~\ref{def:mpdo_twisted_dimer_tensor} has a vertical
    decomposition in the sense of
    Definition~\ref{def:vertical_cf_mpdo}, with the two rescaled flag
    sectors $\widehat M_0$, $\widehat M_1$ of
    Definition~\ref{def:mpdo_twisted_dimer_flag_sectors}, multiplicity one
    each, and weight $\mu=\frac58$. Let $U$ be the permutation matrix
    reordering the original physical index
    $(l,r,f)\mapsto(f,(l,r))$, so that $U$ maps $\C^8$ onto the direct sum
    of the two flag subspaces, each isomorphic to $\C^4$.
    Writing $\widetilde M_{ab}$ for the matrices in the vertical direction,
    we have $UU^\dagger=U^\dagger U=\Id$ and, for every horizontal bond pair
    $(a,b)$,
    \begin{align}
        U\,\widetilde M_{ab}\,U^\dagger
         & =\mu\,\widehat M_0^{ab}\oplus\mu\,\widehat M_1^{ab},
        \notag                                                              \\
        \widetilde M_{ab}
         & =U^\dagger(\mu\,\widehat M_0^{ab}\oplus\mu\,\widehat M_1^{ab})U.
        \notag
    \end{align}
    Each flag sector is injective at one site and therefore proportional to
    a normal tensor with spectral radius one. At every positive length $N$,
    the matrix product vector of the vertically viewed tensor is
    $\mu^N$ times the sum of those of the two sectors, and the matrix product
    vectors of the two sectors are linearly independent. The fact that these
    sectors already have the normalization required of normal tensors in the
    CPSV sense is verified below.
\end{theorem}
\begin{proof}\leanok
    \uses{def:mpdo_twisted_dimer_flag_sectors, lem:mpdo_twisted_dimer_flag_sector_operators, def:vertical_cf_mpdo}
    The matrix of the twisted dimer between two physical indices with
    different flag values vanishes, because its coefficient contains
    $\delta_{ff'}$. Between two indices with common flag value $f$, it is
    $\mu$ times the matrix of $\widehat M_f$ by
    Definition~\ref{def:mpdo_twisted_dimer_flag_sectors}. Conjugating by the
    permutation matrix therefore places the two sectors on the diagonal
    blocks, and the reconstruction follows from $U^\dagger U=\Id$.
    Traces of products are unchanged by conjugation, so
    $\ket{V^{(N)}(\widetilde M)}
        =\mu^N\bigl(\ket{V^{(N)}(\widehat M_0)}
        +\ket{V^{(N)}(\widehat M_1)}\bigr)$.

    For every matrix unit $E_{(p,q),(p',q')}$ of the two-qubit space, the
    matrix of $\widehat M_f$ at the bond pair $((p,p',0),(q,q',0))$ is
    $C_0[p,p']/(2\mu)$ times that matrix unit, with
    $C_0[p,p']\in\{\frac12,\frac38\}$ nonzero. Hence the matrices of each
    sector span the full matrix algebra, giving one-site injectivity.
    For linear independence, evaluate the closed
    sector operators of Lemma~\ref{lem:mpdo_twisted_dimer_flag_sector_operators}
    on two cyclically bond-matched words: on the constant word of bond index
    $(0,0,0)$ both sectors give $(2/5)^N$, since the one-site weight is
    $C_0[0,0]/(2\mu)=\frac25$; on the word with bond index $(0,0,1)$ at
    the first site and $(0,0,0)$ elsewhere, the sector $f$ gives
    $(\tau_1)_{ff}\frac{3}{10}(2/5)^{N-1}$, since
    $C_1[0,0]/(2\mu)=\frac3{10}$, and the two flag signs are opposite.
    A vanishing linear combination
    $g_0\ket{V^{(N)}(\widehat M_0)}
        +g_1\ket{V^{(N)}(\widehat M_1)}=0$
    thus forces $g_0+g_1=0$ and $g_0-g_1=0$.
\end{proof}

\begin{theorem}[Vertical BNT algebra of the twisted quantum dimer]%
    \label{thm:mpdo_twisted_dimer_bnt_algebra_tensor_clause}
    \lean{MPOTensor.TwistedDimer.conjTranspose_flagMPO_mul_self}
    \lean{MPOTensor.TwistedDimer.star_flagCoef_mul_self}
    \lean{MPOTensor.TwistedDimer.flagFamily_isLeftCanonical}
    \lean{MPOTensor.TwistedDimer.flagFamily_isNormalTensor}
    \lean{MPOTensor.TwistedDimer.flagFamily_isCPSVBasisOfNormalTensors}
    \lean{MPOTensor.TwistedDimer.twoLabelChiTracePowerForm}
    \lean{MPOTensor.TwistedDimer.sectorWeight_hasIdempotentCoefficientForm}
    \lean{MPOTensor.TwistedDimer.flagAlgebraClause}
    \lean{MPOTensor.TwistedDimer.T_bntAlgebraTensorClause}
    \lean{MPOTensor.TwistedDimer.T_hasBNTAlgebraTensorClause}
    \leanok
    \uses{def:mpdo_bnt_algebra_tensor_clause, thm:mpdo_twisted_dimer_vertical_cf, thm:mpdo_twisted_dimer_flag_sector_product_law, thm:mpdo_twisted_dimer_two_label_rescaling_stable}
    The vertical sectors of the twisted quantum dimer of
    Definition~\ref{def:mpdo_twisted_dimer_tensor} satisfy the algebra and
    idempotent conditions of
    Definition~\ref{def:mpdo_bnt_algebra_tensor_clause}, using the
    decomposition of Theorem~\ref{thm:mpdo_twisted_dimer_vertical_cf} and
    the structure coefficients of
    Theorem~\ref{thm:mpdo_twisted_dimer_two_label_rescaling_stable}.
    The two flag sectors are normal tensors in the sense of
    \cite[lines~231--235]{Cirac2016MPDO_arXiv}: each is injective at one site
    and left-canonical,
    \begin{align}
        \sum_{a,b}(\widehat M_f^{ab})^\dagger\widehat M_f^{ab}
         & =\Id,
        \notag
    \end{align}
    so its transfer map has spectral radius one. Together with the spanning
    and linear independence already proved, this gives a basis of normal
    tensors of the vertically viewed tensor, as in
    \cite[Proposition~4.13, lines~945--959]{Cirac2016MPDO_arXiv}.
    The positive diagonal matrices
    $\chi_{ff',f+f'}=\alpha$, $\chi_{ff',f+f'+1}=\beta$ give
    $c^{(L)}_{ff'g}=\tr(\chi_{ff'g}^L)$, and the closed operators satisfy
    the product law of
    Theorem~\ref{thm:mpdo_twisted_dimer_flag_sector_product_law}.
    The multiplicity matrices are $\mu_0=\mu_1=[\mu]$, so
    $m_0=m_1=\tr([\mu])=\frac58$. The vector $m$ is an idempotent for the
    multiplication induced by $c^{(1)}$, since
    \begin{align}
        m_g & =\sum_{f,f'}c^{(1)}_{ff'g}\,m_fm_{f'}
        =2(\alpha+\beta)\mu^2=\frac85\cdot\frac{25}{64}=\frac58.
        \notag
    \end{align}
    No positive rescaling of these two sectors makes the structure
    coefficients length independent.
\end{theorem}
\begin{proof}\leanok
    \uses{thm:mpdo_twisted_dimer_vertical_cf, thm:mpdo_twisted_dimer_flag_sector_product_law, thm:is_normal_tensor_of_is_normal_left_canonical}
    The adjoint of a sector matrix times that matrix is the diagonal matrix
    unit at the right bits $(p',q')$ of the bond pair with the squared modulus
    of the coefficient, and the flag sign squares to one. Summing over the
    bond pairs $((p,p',k),(q,q',k))$ with the right bits fixed gives, at each
    diagonal entry,
    \begin{align}
        \frac{1}{(2\mu)^2}\sum_{k}\sum_{p,q}C_k[p,p']^2
         & =\frac{2\cdot2\cdot(\tfrac14+\tfrac9{64})}{4\cdot\tfrac{25}{64}}
        =1,
        \notag
    \end{align}
    and the off-diagonal entries vanish. Together with one-site injectivity
    from Theorem~\ref{thm:mpdo_twisted_dimer_vertical_cf}, this makes each
    sector a normal tensor with spectral radius one by
    Theorem~\ref{thm:is_normal_tensor_of_is_normal_left_canonical}.
    The spanning and linear independence are those of
    Theorem~\ref{thm:mpdo_twisted_dimer_vertical_cf}. The product law is
    Theorem~\ref{thm:mpdo_twisted_dimer_flag_sector_product_law}.
    For each output label $g$, two pairs $(f,f')$ contribute $\alpha$ and
    two contribute $\beta$, so
    $\sum_{f,f'}c^{(1)}_{ff'g}m_fm_{f'}=2(\alpha+\beta)\mu^2=\mu=m_g$.
\end{proof}

\begin{lemma}[Doubled-index matrices and their rank-one factorization]
    \label{lem:mpdo_bnt_label_rescaling_stable_doubled_letter_factorization}
    \lean{MPOTensor.RescalingStableLengthDependentRFP.toMPSTensor_R_apply}
    \lean{MPOTensor.RescalingStableLengthDependentRFP.R_eq_vecMulVec}
    \leanok
    \uses{thm:mpdo_bnt_label_rescaling_stable_length_dependence}
    Identify each physical index $p\in\{0,1,2,3\}$ with a pair
    $(p_1,p_2)\in\{0,1\}^2$, and similarly write $q=(q_1,q_2)$. The
    doubled-index matrix at the paired index $(p,q)$ is the MPO matrix $R^{pq}$.
    This bond matrix has the outer-product factorization
    \begin{align}
        R^{pq}_{ab}
         & =\frac{25}{32}A^a_{p_1q_1}A^b_{p_2q_2},
        \notag
    \end{align}
    or equivalently $R^{pq}=\frac{25}{32}u_{pq}v_{pq}^{\mathsf T}$, where
    $u_{pq}(a)=A^a_{p_1q_1}$ and $v_{pq}(b)=A^b_{p_2q_2}$.
\end{lemma}

\begin{proof}\leanok
    The doubled-index reading gives $R^{pq}$ by definition. Expanding the
    Kronecker-product definition of $R$ and using that the entries of the
    four matrices $A^a$ are real gives
    $R^{pq}=\frac{25}{32}u_{pq}v_{pq}^{\mathsf T}$.
\end{proof}

\begin{lemma}[Rank-one form of the dimer transfer map]
    \label{lem:mpdo_bnt_label_rescaling_stable_transfer_rank_one}
    \lean{MPOTensor.RescalingStableLengthDependentRFP.%
        transferMap_R_toMPSTensor_eq}
    \lean{MPOTensor.RescalingStableLengthDependentRFP.cFun_fixedDiag}
    \leanok
    \uses{thm:mpdo_bnt_label_rescaling_stable_length_dependence}
    Set $s=(\frac45,\frac45,\frac35,\frac35)$,
    $F=\diag(s_0^2,s_1^2,s_2^2,s_3^2)$, and
    \begin{align}
        \lambda(X) & =\left(\frac{25}{32}\right)^2
        \sum_{b=0}^3s_b^2X_{bb}.
        \notag
    \end{align}
    Then the transfer map of $R$ has rank-one form
    \begin{align}
        \E_R(X)    & =\lambda(X)F,
        \notag                         \\
        \lambda(F) & =\frac{337}{512}.
        \notag
    \end{align}
\end{lemma}

\begin{proof}\leanok
    \uses{lem:mpdo_bnt_label_rescaling_stable_doubled_letter_factorization}
    By
    Lemma~\ref{lem:mpdo_bnt_label_rescaling_stable_doubled_letter_factorization},
    the doubled-index matrices are
    $R^{pq}=\frac{25}{32}u_{pq}v_{pq}^{\mathsf T}$.
    Substituting these factors into the transfer-map sum and summing over
    the physical indices gives $\E_R(X)=\lambda(X)F$.
    Substituting $X=F$ and evaluating the four diagonal entries gives
    $\lambda(F)=337/512$.
\end{proof}

\begin{lemma}[One-site transfer quasi-idempotence of the dimer tensor]
    \label{lem:mpdo_bnt_label_rescaling_stable_transfer_quasi_idempotent}
    \lean{MPOTensor.RescalingStableLengthDependentRFP.%
        transferMap_R_toMPSTensor_idem}
    \leanok
    \uses{thm:mpdo_bnt_label_rescaling_stable_length_dependence}
    The transfer map of the dimer tensor $R$ satisfies
    $\E_R^2=\frac{337}{512}\E_R$.
\end{lemma}

\begin{proof}\leanok
    \uses{lem:mpdo_bnt_label_rescaling_stable_transfer_rank_one}
    By Lemma~\ref{lem:mpdo_bnt_label_rescaling_stable_transfer_rank_one},
    \begin{align}
        \E_R(\E_R(X))
         & =\lambda(X)\lambda(F)F
        =\frac{337}{512}\E_R(X).
        \notag
    \end{align}
\end{proof}
